\begin{abstract}
\noindent
Une campagne antérieure d'optimisation FO-LQI $\to$ FO-PIDF pour un
convertisseur Buck non idéal de 48~V avait consommé plus de 1800
évaluations sans atteindre la faisabilité robuste, et attribuait cet échec à
l'écrêtage du gain proportionnel sur une borne de paramétrisation. Nous
proposons un \emph{audit de saturation des bornes}, qui distingue saturations
structurelles et pathologiques, et nous l'appliquons à une chaîne
reconstruite dont la projection reproduit la proposition d'identité exacte à
\ChiffreErreurPropUn{} près. Sur \ChiffreAuditN{} tirages de la boîte de
décision, $\log_{10}k_p$ ne touche jamais sa borne : l'écrêtage n'explique
pas l'échec. Une campagne confirmatoire préenregistrée (protocole figé par
empreintes, trois bras, six métaheuristiques, trente graines appariées,
540 cellules) sur un simulateur commuté qui reproduit la conduction
discontinue ne produit \emph{aucun} correcteur robustement faisable : 0/30
graines par méthode et par bras, intervalle de Wilson $[0;\ChiffreWilsonHaut]$.
La paramétrisation conforme au cadrage n'atteint même pas la faisabilité
nominale ; libérer la pondération de consigne~$b$ la rend accessible dans
\ChiffreNomFaisCinqB{} cellules sur 180, sans robustesse. Un diagnostic
physique identifie un verrou indépendant du correcteur : à rapport cyclique
minimal, le convertisseur en conduction discontinue se stabilise sous le
seuil de tension minimale dans \ChiffrePlancherN{} des \ChiffrePlancherTot{}
cas d'incertitude. Nous montrons enfin que des correcteurs faisables au sens
des six contraintes peuvent entretenir un cycle limite que ces contraintes ne
voient pas. Ces résultats portent sur un modèle simulé et un ensemble fini de
cas ; ils ne constituent ni un certificat de robustesse ni une preuve
d'infaisabilité.
\end{abstract}

\noindent\textbf{Mots-clés :} convertisseur Buck ; commande d'ordre
fractionnaire ; FO-LQI ; FO-PIDF ; métaheuristiques ; audit de paramétrisation ;
conduction discontinue ; reproductibilité.
